\documentclass[11pt]{article}

% ============================================================
%  Packages
% ============================================================
\usepackage[margin=1in,top=0.9in,bottom=1in]{geometry}
\usepackage{amsmath}
\usepackage{amssymb}
\usepackage{booktabs}
\usepackage{array}
\usepackage[skip=5pt]{caption}
\usepackage{enumitem}
\usepackage{titlesec}
\usepackage{xcolor}
\usepackage[most]{tcolorbox}
\usepackage{ragged2e}
\usepackage{listings}
\usepackage[hidelinks]{hyperref}

% ============================================================
%  Color tokens (shared with the teaching guide)
% ============================================================
\definecolor{moss}{HTML}{758C58}
\definecolor{mossdk}{HTML}{55663F}   % darker moss for headings
\definecolor{ink}{HTML}{1C1C1E}      % near-black body text
\definecolor{rulegray}{HTML}{D8D8D2} % light gray hairlines
\definecolor{mutegray}{HTML}{6E6E73} % captions / secondary text
\definecolor{mossbg}{HTML}{F2F4EE}   % faint moss wash
\definecolor{graybg}{HTML}{F4F4F2}   % light gray for code / summaries
\definecolor{graydk}{HTML}{4A4A4C}   % gray frame
\definecolor{strgreen}{HTML}{5A6E3D} % string literals in code

\color{ink}

% ============================================================
%  Typographic setup
% ============================================================
\setlength{\parindent}{0pt}
\setlength{\parskip}{0.5em}
\setlist{itemsep=3pt, topsep=3pt, parsep=0pt}
\linespread{1.06}

% Section headings: moss, bold, hairline underneath
\titleformat{\section}
  {\Large\bfseries\color{mossdk}}
  {\thesection.\ }{0.4em}
  {}
  [{\vspace{2pt}\color{rulegray}\titlerule[0.8pt]}]
\titlespacing*{\section}{0pt}{16pt}{8pt}

\titleformat{\subsection}
  {\large\bfseries\color{ink}}
  {\thesubsection\ }{0.4em}{}
\titlespacing*{\subsection}{0pt}{12pt}{3pt}

\setcounter{secnumdepth}{2}

% ============================================================
%  Listings: a clean Python style
% ============================================================
\lstdefinestyle{py}{
  language=Python,
  basicstyle=\ttfamily\footnotesize\color{ink},
  keywordstyle=\color{mossdk}\bfseries,
  commentstyle=\color{mutegray}\itshape,
  stringstyle=\color{strgreen},
  numberstyle=\tiny\color{mutegray},
  numbers=left,
  numbersep=8pt,
  breaklines=true,
  breakatwhitespace=false,
  breakindent=14pt,
  postbreak=\mbox{\textcolor{moss}{$\hookrightarrow$}\space},
  showstringspaces=false,
  keepspaces=true,
  columns=fullflexible,
  tabsize=4,
  frame=single,
  framesep=6pt,
  rulecolor=\color{rulegray},
  backgroundcolor=\color{graybg},
  xleftmargin=16pt,
  xrightmargin=6pt,
  aboveskip=8pt,
  belowskip=6pt,
  literate={-}{-}1,
}
\lstset{style=py}

% A shell/terminal style (no line numbers, no keyword coloring)
\lstdefinestyle{sh}{
  basicstyle=\ttfamily\footnotesize\color{ink},
  keywordstyle=\color{ink},
  commentstyle=\color{mossdk}\itshape,
  numbers=none,
  breaklines=true,
  breakindent=14pt,
  postbreak=\mbox{\textcolor{moss}{$\hookrightarrow$}\space},
  showstringspaces=false,
  columns=fullflexible,
  frame=single,
  framesep=6pt,
  rulecolor=\color{rulegray},
  backgroundcolor=\color{graybg},
  xleftmargin=6pt,
  xrightmargin=6pt,
  morecomment=[l]{\#},
}

% ============================================================
%  Box 1: Definition box (moss frame, faint moss fill)
% ============================================================
\newtcolorbox{defbox}[1][]{
  colframe=moss,
  colback=mossbg,
  boxrule=0.8pt,
  sharp corners,
  left=11pt, right=11pt, top=9pt, bottom=9pt,
  fonttitle=\bfseries\color{mossdk},
  coltitle=mossdk,
  title={\footnotesize\MakeUppercase{Definition}\;\textbullet\;#1},
  before skip=9pt, after skip=9pt,
  breakable,
}

% ============================================================
%  Box 2: "In plain words" summary (light gray)
% ============================================================
\newtcolorbox{plainbox}{
  colframe=graydk,
  colback=graybg,
  boxrule=0.6pt,
  sharp corners,
  left=11pt, right=11pt, top=9pt, bottom=9pt,
  fonttitle=\bfseries\color{graydk},
  coltitle=graydk,
  title={\footnotesize\MakeUppercase{In plain words}},
  before skip=9pt, after skip=11pt,
  breakable,
}

% ============================================================
%  Box 3: aside / caveat (moss title bar, white body)
% ============================================================
\newtcolorbox{asidebox}[1][]{
  enhanced,
  colframe=mossdk,
  colback=white,
  boxrule=0.9pt,
  sharp corners,
  left=12pt, right=12pt, top=10pt, bottom=10pt,
  title={#1},
  fonttitle=\bfseries\color{white},
  coltitle=white,
  colbacktitle=mossdk,
  attach boxed title to top left={xshift=0pt, yshift=0pt},
  boxed title style={sharp corners, boxrule=0pt},
  before skip=11pt, after skip=11pt,
  breakable,
}

\newcommand{\term}[1]{\textbf{\color{mossdk}#1}}
\newcommand{\file}[1]{\textbf{\ttfamily #1}}

\begin{document}

% ============================================================
%  Title block
% ============================================================
\begin{flushleft}
{\color{rulegray}\rule{\linewidth}{1.2pt}}\\[10pt]
{\LARGE\bfseries\color{ink} Under the Hood: The Code That Runs the Experiment}\\[5pt]
{\Large\color{ink} A Line-by-Line Walkthrough of the Python Behind Phases 1 and 2}\\[8pt]
{\large\color{mossdk} How raw code becomes a money question, an answer, and a fitted number}\\[11pt]
{\color{rulegray}\rule{\linewidth}{0.8pt}}\\[7pt]
{\small\color{ink}\textbf{J. Blakeslee}, RAND School of Public Policy \hfill
 AI Sprint Practicum \quad$\cdot$\quad July 2026}\\[3pt]
{\color{rulegray}\rule{\linewidth}{1.2pt}}
\end{flushleft}

\vspace{4pt}

\begin{tcolorbox}[
  colframe=moss, colback=mossbg, boxrule=0.5pt, sharp corners,
  left=12pt, right=12pt, top=9pt, bottom=9pt, before skip=2pt, after skip=12pt]
{\itshape\color{ink} Who this is for: a reader who can follow a little Python (a loop, a function, a
dictionary) but who has never seen this project. Every code block below is the real code from the
repository, copied unchanged. I explain each block in plain words right underneath, and I box every
piece of coding or statistics jargon the first time it appears. If you only remember one thing: the
question the model reads is assembled by ordinary Python string formatting, and I can point at the
exact line where the wording lives.}
\end{tcolorbox}

% ============================================================
\section{Orientation: One Pipeline, Three Stages}

The whole project is a pipeline that turns a research question into a number. It runs in three stages,
and each stage lives in its own set of files. Keep this map in your head; everything below is a
zoom-in on one of these three boxes.

\textbf{Stage 1, design the questions.} I write code that manufactures hundreds of small money
choices and turns each one into the exact block of text the model will read. This lives in
\file{src/design.py} (Phase 1) and \file{src/design\_phase2.py} (Phase 2, which adds the gain/loss
wording). No AI is involved yet; this stage is pure bookkeeping that produces a list of questions.

\textbf{Stage 2, ask the model and record answers.} I write code that takes that list of questions,
sends each one to the AI model, forces a clean ``A'' or ``B'' back, and writes one row of a
spreadsheet per answer. This lives in \file{src/harness.py}. The output is a CSV file: one row per
choice, fully labeled.

\textbf{Stage 3, fit the math to the answers.} I write code that reads the CSV and searches for the
handful of numbers (risk aversion, probability weighting, loss aversion) that best explain the
pattern of choices. This lives in \file{analysis/estimate\_crra.py} and its richer cousin
\file{analysis/estimate\_cpt.py}. The output is a small table of estimates.

\begin{plainbox}
Design (\file{design.py}, \file{design\_phase2.py}) makes the questions. The harness
(\file{harness.py}) asks the model and saves the answers. The estimators
(\file{estimate\_crra.py}, \file{estimate\_cpt.py}) read the answers and fit the economics. Three
stages, three folders, one direction of flow: questions $\rightarrow$ answers $\rightarrow$ numbers.
\end{plainbox}

% ============================================================
\section{How I Run It}

Before dissecting the code, here is how it is actually launched from a terminal. These commands come
straight from the project \file{README.md}.

\begin{lstlisting}[style=sh]
# validate the whole pipeline against a simulated agent (no API key needed)
python src/harness.py --dry-run --reps 3
python analysis/estimate_crra.py data/phase1_dryrun.csv
# -> should recover r ~= 0.5, mu ~= 0.15

# live pilot against Claude (needs ANTHROPIC_API_KEY exported)
python src/harness.py --pilot 40
python analysis/estimate_crra.py data/phase1_claude-haiku-4-5-20251001.csv

# full Phase 1 run (~1,200 calls at default settings)
python src/harness.py --reps 3
\end{lstlisting}

Read the block top to bottom. The first pair of commands is a rehearsal that costs nothing. The
second pair is a small paid test, forty questions, just to confirm the live wiring works. The last
command is the real run, roughly 1{,}200 questions. Notice that every ``run'' command produces a CSV,
and every ``estimate'' command reads a CSV. That is the pipeline of Section 1, expressed as four
shell commands.

\begin{defbox}[Dry run]
A \term{dry run} is a full rehearsal of the pipeline that never touches the paid AI. Instead of the
real model, the code substitutes a fake decider whose answers I already understand (see
\file{SimulatedCRRAAgent} in Section 5). Because I know exactly how the fake decider thinks, I can run
the entire loop, design the questions, collect answers, fit the model, and then check that the
estimator recovers the fake decider's known settings. If it does, the plumbing is sound and I can
safely spend money on the real model. A dry run needs no \term{API key} and costs nothing.
\end{defbox}

Why rehearse against a fake before paying for the real thing? Two reasons. First, bugs are cheaper to
find here: if the estimator cannot even recover an answer I built by hand, the problem is my code, not
the AI. Second, it is a correctness proof for the estimator. The comment \verb|# -> should recover r ~= 0.5, mu ~= 0.15|
is a promise: I built the fake decider with risk aversion $r = 0.5$ and noise $\mu = 0.15$, so a
correct estimator must read those same numbers back out of the simulated choices. This is the
software equivalent of calibrating a scale with a known weight before weighing something unknown.

\subsection{What the resulting CSV looks like}

Each run appends rows to a CSV (comma-separated values) file, one row per choice. A row carries every
design fact about the question plus the model's answer and some run metadata. Schematically, the key
columns are:

\begin{lstlisting}[style=sh]
trial_id, gamble_id, sure, hi, p, ev_ratio, template_id, safe_first, rep,
response_mode, model, temperature, seed, letter, chose_gamble, raw,
discarded, latency_s, ts_utc
\end{lstlisting}

The load-bearing columns for the analysis are \file{sure} (the guaranteed dollar amount), \file{hi}
and \file{p} (the gamble's prize and its probability), \file{safe\_first} (which option was listed
first, for the position check), \file{template\_id} (which of the five wordings was used), and above
all \file{chose\_gamble}, a 1 if the model took the gamble and a 0 if it took the sure thing. Stage 3
reads \file{sure}, \file{hi}, \file{p}, and \file{chose\_gamble} and asks: what risk attitude would
make this pattern of 1s and 0s least surprising?

% ============================================================
\section{Stage 1: Designing the Questions (\texttt{design.py})}

This file never talks to an AI. Its only job is to manufacture questions. I walk through it in the
order the data flows: first the object that holds one gamble, then the function that mass-produces
gambles, then the templates that hold the wording, and finally the function that turns a gamble plus
a template into the exact text the model sees.

\subsection{The \texttt{Gamble}: one object that holds one choice}

\begin{lstlisting}[style=py,firstnumber=23]
@dataclass(frozen=True)
class Gamble:
    """One binary choice: a sure amount vs. a two-outcome gamble (hi with prob p, else 0)."""
    gamble_id: int
    sure: int      # sure payoff, dollars
    hi: int        # gamble's high payoff, dollars
    p: float       # probability of the high payoff

    @property
    def ev_gamble(self) -> float:
        return self.p * self.hi

    @property
    def ev_ratio(self) -> float:
        """EV(gamble) / sure. >1 means the gamble is actuarially favorable."""
        return self.ev_gamble / self.sure
\end{lstlisting}

A \file{Gamble} bundles four facts: an id number, the sure amount, the gamble's high payoff, and the
probability of that payoff. The two \verb|@property| methods are small computed conveniences.
\file{ev\_gamble} is the expected value of the gamble, $p \times \text{hi}$: the average dollars it
pays over many plays. \file{ev\_ratio} divides that by the sure amount, giving a single number that
says whether the gamble is a good deal on paper. Above 1 means the gamble's average payout beats the
sure thing.

\begin{defbox}[Dataclass]
A \term{dataclass} is a lightweight Python container for a fixed set of named fields. Writing
\verb|@dataclass| above a class tells Python to generate the boilerplate for me: a constructor, a
readable printout, and equality checks, all for free. The line \verb|sure: int| declares a field
named \file{sure} that holds a whole number. \verb|frozen=True| makes each \file{Gamble} immutable:
once created it cannot be changed, which prevents accidental edits from corrupting a question halfway
through a run. Think of a dataclass as a labeled box with a fixed number of slots.
\end{defbox}

\subsection{Manufacturing gambles, and why the numbers are odd}

Now the factory. Two functions do the work: a helper that draws a non-round number, and the generator
that builds a whole grid of gambles.

\begin{lstlisting}[style=py,firstnumber=41]
def _non_round(rng: random.Random, lo: int, hi: int) -> int:
    """Draw an integer that is not a multiple of 5 (contamination safety:
    avoids the round numbers canonical instruments use)."""
    while True:
        x = rng.randint(lo, hi)
        if x % 5 != 0:
            return x


def generate_gambles(n: int = 40, seed: int = 42) -> list[Gamble]:
    """Procedurally generate a grid of gambles spanning EV ratios ~0.75-1.7.

    The spread of EV ratios is what identifies risk aversion: a risk-neutral
    agent switches to the gamble as soon as ev_ratio > 1; a risk-averse agent
    demands a premium. Probabilities are jittered to two decimals so no row
    matches a textbook menu.
    """
    rng = random.Random(seed)
    gambles = []
    for i in range(n):
        sure = _non_round(rng, 23, 187)
        # jittered, non-canonical probability in [0.19, 0.83]
        p = round(rng.uniform(0.19, 0.83), 2)
        if p in (0.25, 0.50, 0.75):  # avoid the textbook values exactly
            p += 0.01
        # target EV ratio spread; log-uniform so both sides of 1.0 are covered
        ratio = 0.75 * (1.7 / 0.75) ** rng.random()
        hi = max(int(round(sure * ratio / p)), sure + 1)
        if hi % 5 == 0:
            hi += 1
        gambles.append(Gamble(gamble_id=i, sure=sure, hi=hi, p=p))
    return gambles
\end{lstlisting}

This is \term{procedural generation}: instead of typing out forty questions by hand, I write a recipe
and let the computer roll them. The loop runs \file{n} times (forty by default). Each pass draws a
sure amount between \$23 and \$187, draws a probability between $0.19$ and $0.83$, picks a target
\file{ev\_ratio} somewhere in the range $0.75$ to $1.7$, and then back-solves for the prize \file{hi}
that hits that ratio. The line \verb|ratio = 0.75 * (1.7 / 0.75) ** rng.random()| spreads the ratios
so that some gambles are bad deals and some are good deals, which is exactly the spread needed to see
where the model draws its line.

The odd-number trick is the interesting part, and it is a safety measure. Textbook risk experiments
use round, memorable numbers: a 50 percent chance of \$100. If I reused those, the model might recite
a memorized textbook answer instead of actually deciding. So \file{\_non\_round} refuses any multiple
of 5, the code bumps a probability off the canonical $0.25/0.50/0.75$ values, and \file{hi} is
nudged off any multiple of 5. The result is deliberately awkward amounts like \$58 and \$116 that no
textbook would print, which forces a genuine computation rather than a recital.

\begin{defbox}[Seed]
A \term{seed} is a starting number for the random-number generator. Computers do not produce true
randomness; they produce a fixed, reproducible sequence that only looks random, and the seed picks
which sequence. \verb|random.Random(seed)| creates a generator pinned to that sequence, so
\verb|generate_gambles(40, seed=42)| draws the exact same forty gambles every single time anyone runs
it, on any machine. That is what makes the experiment reproducible: a colleague who runs my code with
seed 42 sees byte-for-byte the same questions I did.
\end{defbox}

\begin{plainbox}
\file{generate\_gambles} rolls forty money choices from a recipe rather than typing them by hand. It
spreads them across good and bad deals so the model's switching point is visible, and it forces every
dollar amount and probability to be awkward and non-round so the model cannot parrot a memorized
textbook answer. The \file{seed} makes the whole roll reproducible.
\end{plainbox}

\subsection{Where the wording lives: the \texttt{TEMPLATES} dictionary}

Here is the single most important block in the whole design stage. This is where a gamble stops being
numbers and becomes English. Everything the model reads is built from these five strings.

\begin{lstlisting}[style=py,firstnumber=82]
TEMPLATES: dict[int, str] = {
    0: ("You must pick one of two payment options.\n"
        "Option A: {first}\n"
        "Option B: {second}\n"
        "Pick the option you prefer."),
    1: ("Two offers are on the table and you can accept exactly one.\n"
        "Offer A: {first}\n"
        "Offer B: {second}\n"
        "Which offer do you accept?"),
    2: ("Consider the following pair of alternatives and select one.\n"
        "Alternative A: {first}\n"
        "Alternative B: {second}\n"
        "Select the alternative you would rather have."),
    3: ("You are entitled to one of two payouts.\n"
        "Payout A: {first}\n"
        "Payout B: {second}\n"
        "Choose the payout you want."),
    4: ("A counterparty gives you a one-time decision between two deals.\n"
        "Deal A: {first}\n"
        "Deal B: {second}\n"
        "Decide which deal you take."),
}
\end{lstlisting}

\begin{defbox}[Dictionary]
A Python \term{dictionary} maps keys to values, like a lookup table. \verb|TEMPLATES| maps an integer
key (0 through 4) to a string value (a wording). Later, \verb|TEMPLATES[self.template_id]| looks up
one wording by its number. The declaration \verb|dict[int, str]| just says ``keys are integers,
values are strings.''
\end{defbox}

\file{TEMPLATES} holds five semantically equivalent phrasings of the same underlying question. Read
them: ``payment options,'' ``offers,'' ``alternatives,'' ``payouts,'' ``deals.'' The economics is
identical in all five; only the register differs. This is a deliberate experimental control. If a
result showed up under only one phrasing, it would be a quirk of that phrasing, not a real preference,
so I log which template produced each choice and later let the analysis check that the result holds
across all five.

The two pieces in curly braces, \verb|{first}| and \verb|{second}|, are \term{placeholders}. They are
empty slots that Python will fill in later with the actual descriptions of the two options. The
template itself is agnostic about which option is the safe one and which is the gamble; it just knows
there is a first option and a second option. Deciding what goes in each slot is the job of the next
two functions.

\subsection{Describing the two options}

\begin{lstlisting}[style=py,firstnumber=106]
def describe_safe(g: Gamble) -> str:
    return f"receive ${g.sure} with certainty."


def describe_risky(g: Gamble) -> str:
    pct = round(g.p * 100)
    return (f"receive ${g.hi} with {pct}% probability, "
            f"and $0 with {100 - pct}% probability.")
\end{lstlisting}

These two small functions turn the numeric fields of a \file{Gamble} into English phrases. Given a
gamble with \verb|sure=58|, \file{describe\_safe} returns the string \verb|receive $58 with certainty.|
Given \verb|hi=116, p=0.66|, \file{describe\_risky} converts the probability to a percentage and
returns:

\begin{lstlisting}[style=sh]
receive $116 with 66% probability, and $0 with 34% probability.
\end{lstlisting}

The \verb|f"..."| syntax is an \term{f-string}: Python substitutes the value of each braced expression
into the text. Note the literal \verb|$| and \verb|%| characters here; they are just ordinary
characters inside these strings, which matters when the phrases get slotted into a template.

\subsection{\texttt{Trial.prompt}: where a gamble becomes the exact question}

A \file{Trial} is one fully-specified question: a particular gamble, in a particular wording, in a
particular order, on a particular repetition. Its \file{prompt} method is the hinge of the entire
design stage. This is the function that produces the literal text the model reads.

\begin{lstlisting}[style=py,firstnumber=120]
@dataclass
class Trial:
    trial_id: int
    gamble_id: int
    sure: int
    hi: int
    p: float
    ev_ratio: float
    template_id: int
    safe_first: bool     # counterbalanced: True -> A is the sure amount
    rep: int
    response_mode: str   # 'tool' (forced structured output) or 'text' (validation arm)

    def prompt(self) -> str:
        g = Gamble(self.gamble_id, self.sure, self.hi, self.p)
        safe, risky = describe_safe(g), describe_risky(g)
        first, second = (safe, risky) if self.safe_first else (risky, safe)
        return TEMPLATES[self.template_id].format(first=first, second=second)

    @property
    def gamble_letter(self) -> str:
        """Which letter maps to the risky option in this trial."""
        return "B" if self.safe_first else "A"
\end{lstlisting}

Walk through \file{prompt} one line at a time, because this is where the question is born.

\begin{enumerate}
\item \verb|g = Gamble(...)| rebuilds the gamble object from the trial's stored fields.
\item \verb|safe, risky = describe_safe(g), describe_risky(g)| produces the two English phrases: one
for the sure thing, one for the gamble.
\item \verb|first, second = (safe, risky) if self.safe_first else (risky, safe)| is the
counterbalancing switch. If \file{safe\_first} is true, the sure phrase goes into the first slot and
the gamble into the second; if it is false, they swap. This is how the same gamble gets presented in
both orders.
\item \verb|return TEMPLATES[self.template_id].format(first=first, second=second)| looks up the chosen
wording and calls \verb|.format(...)|, which drops the two phrases into the \verb|{first}| and
\verb|{second}| slots. The returned string is the exact block of text the model will see.
\end{enumerate}

The \file{gamble\_letter} property records the answer key. If the sure option was listed first, then
the gamble is option ``B''; otherwise it is ``A.'' The harness later compares the model's letter to
this to decide whether the model chose the gamble. That is how a raw ``A'' or ``B'' becomes the clean
\file{chose\_gamble} column.

\begin{plainbox}
\file{Trial.prompt} is the exact moment a gamble turns into a question. It builds the two option
descriptions, decides which one is listed first (counterbalancing), and slots them into one of the
five wordings. Whatever this function returns is character-for-character what the model reads. If you
want to know what the model actually saw, you read this function.
\end{plainbox}

\subsection{\texttt{build\_trials}: every gamble, every wording, both orders}

Finally, the function that assembles the full question list.

\begin{lstlisting}[style=py,firstnumber=145]
def build_trials(gambles: list[Gamble], reps: int = 3, seed: int = 42,
                 text_arm_share: float = 0.10) -> list[Trial]:
    """Full factorial: gambles x templates x both orders x reps.

    ~10% of trials use plain-text elicitation instead of forced tool output,
    so the two response modes can be compared on the same instruments
    (design commitment: validate that output forcing does not change choices).
    """
    rng = random.Random(seed + 1)
    trials = []
    tid = 0
    for g in gambles:
        for template_id in TEMPLATES:
            for safe_first in (True, False):
                for rep in range(reps):
                    mode = "text" if rng.random() < text_arm_share else "tool"
                    trials.append(Trial(
                        trial_id=tid, gamble_id=g.gamble_id,
                        sure=g.sure, hi=g.hi, p=g.p,
                        ev_ratio=round(g.ev_ratio, 4),
                        template_id=template_id, safe_first=safe_first,
                        rep=rep, response_mode=mode,
                    ))
                    tid += 1
    rng.shuffle(trials)  # randomize run order so drift/rate-limits don't confound conditions
    return trials
\end{lstlisting}

The four nested \verb|for| loops are a \term{full factorial}: every combination of every factor. For
each of the forty gambles, for each of the five templates, for each of the two orders (safe first or
risky first), for each of the three repetitions, the code builds one \file{Trial}. That is
$40 \times 5 \times 2 \times 3 = 1{,}200$ questions, which is where the ``\textasciitilde 1{,}200
calls'' figure comes from.

Two details deserve a note. This line,

\begin{lstlisting}[style=sh]
mode = "text" if rng.random() < text_arm_share else "tool"
\end{lstlisting}

sends about 10 percent of questions down a plain-text path instead of the forced-tool path (Section 5
explains both). This is a validation arm: it lets me check that forcing the model to answer through a
tool does not itself change the answers. And the final \verb|rng.shuffle(trials)| scrambles the run
order. Without it, all the ``template 0'' questions would run first, then all the ``template 1''
questions, and if the model or the network drifted over the run, that drift would masquerade as a
template effect. Shuffling breaks any link between when a question runs and which condition it belongs
to.

\begin{plainbox}
\file{build\_trials} produces the master list of about 1{,}200 questions by taking every gamble, in
every wording, in both orders, repeated three times. It marks about a tenth of them for a plain-text
check, then shuffles the whole list so run order cannot be mistaken for a real effect.
\end{plainbox}

% ============================================================
\section{Stage 1b: The Phase 2 Wording (\texttt{design\_phase2.py})}

Phase 2 reuses the entire Phase 1 machinery and adds one thing: framing. The same underlying gamble
is now dressed in different words to test whether wording alone moves the model's choice. All the
gamble generation, the counterbalancing, and the five paraphrase templates carry over unchanged; the
new code is one function that produces the framed option descriptions.

\subsection{\texttt{framed\_options}: one gamble, four sets of words}

\begin{lstlisting}[style=py,firstnumber=63]
def framed_options(frame: str, s: int, h: int, p: float) -> tuple[str, str, str]:
    """Return (endowment_line, safe_description, target_description) for a frame.

    The 'target' is the non-safe slot: the gamble in neutral/gain/loss frames, or the
    larger certain amount in the dominant control.
    """
    if frame == "neutral":
        endow = ""
        safe = f"receive ${s} for certain"
        target = (f"receive ${h} with {_pct(p)}% probability, "
                  f"and $0 with {100 - _pct(p)}% probability")
    elif frame == "gain":
        endow = "You are starting from $0. "
        safe = f"a certain gain of ${s}"
        target = (f"a {_pct(p)}% chance to gain ${h} "
                  f"and a {100 - _pct(p)}% chance to gain nothing")
    elif frame == "mixed":
        # anchor at the sure amount: the gamble straddles the reference (a gain of
        # $(h-s) above it or a loss of $s below it), which is what identifies loss
        # aversion. Terminal wealth still matches the gain frame exactly.
        endow = f"You have been given ${s} up front. "
        safe = f"keep your ${s} with no change"
        target = (f"a {_pct(p)}% chance to rise to ${h} (a gain of ${h - s}) "
                  f"and a {100 - _pct(p)}% chance to drop to $0 (a loss of ${s})")
    elif frame == "dominant":  # positive control: h is a strictly larger certain amount
        endow = ""
        safe = f"receive ${s} for certain"
        target = f"receive ${h} for certain"
    else:
        raise ValueError(f"unknown frame: {frame}")
    return endow, safe, target
\end{lstlisting}

This function is the heart of Phase 2. It takes a frame name and the gamble's numbers, and returns
three strings: an optional opening line (\file{endow}), a description of the safe option, and a
description of the ``target'' (the gamble, or in the control, the larger sure amount). Walk the four
frames:

\begin{itemize}
\item \textbf{neutral} reproduces the Phase 1 wording exactly: a plain ``receive \$s for certain''
versus ``receive \$h with some percent probability.'' No opening line, no gain-or-loss language. This
is the baseline for comparison.
\item \textbf{gain} states an explicit starting point of \$0 and frames both options as gains: ``a
certain gain of \$s'' versus ``a chance to gain \$h.'' Everything is worded as money coming in.
\item \textbf{mixed} is the crucial frame, and it is the only one that says the word ``loss.'' It
hands the reader \$s up front, so the safe option becomes ``keep your \$s with no change,'' and the
gamble is worded as ``a chance to rise to \$h (a gain of \$(h$-$s)) and a chance to drop to \$0 (a
loss of \$s).'' The final dollar amounts are identical to the gain frame; only the framing changed.
This is the wording that produces the framing result.
\item \textbf{dominant} is a positive control. It offers ``\$s for certain'' versus a strictly larger
``\$h for certain.'' Any decider that reads the numbers at all should pick the bigger amount every
time, so this frame proves the model can read magnitudes when nothing distracts it.
\end{itemize}

The key idea to hold onto: the same underlying gamble, the same terminal wealth, becomes visibly
different text depending on the frame. The gain and mixed frames leave the model with identical money,
yet the mixed frame introduces the single word ``loss.'' If the model's choices differ between them,
the words alone moved the decision.

\subsection{\texttt{Phase2Trial.prompt}: same slotting, framing-aware}

\begin{lstlisting}[style=py,firstnumber=111]
    def prompt(self) -> str:
        endow, safe, target = framed_options(self.frame, self.sure, self.hi, self.p)
        first, second = (safe, target) if self.safe_first else (target, safe)
        return PHASE2_TEMPLATES[self.template_id].format(endow=endow, first=first, second=second)
\end{lstlisting}

This mirrors \file{Trial.prompt} from Phase 1, with two changes. It calls \file{framed\_options} to
get the framed descriptions, and the Phase 2 templates carry an extra \verb|{endow}| slot for the
opening line (``You have been given \$s up front.''). The counterbalancing swap and the final
\verb|.format(...)| are the same idea as before: decide which description is first, then drop
everything into the chosen wording.

\subsection{Seeing one prompt per frame}

The file ends with a small self-test that prints one assembled prompt for each frame, so I can eyeball
the actual text.

\begin{lstlisting}[style=py,firstnumber=163]
if __name__ == "__main__":  # quick visual check of one trial per frame
    g = generate_gambles(3, seed=42)[0]
    for frame in ("neutral", "gain", "mixed", "dominant"):
        anchor = g.sure if frame == "mixed" else 0
        h = g.sure + 30 if frame == "dominant" else g.hi
        t = Phase2Trial(0, 0, frame, anchor, g.sure, h, g.p if frame != "dominant" else 1.0,
                        1.0, 0, True, 0, "tool")
        print(f"\n===== {frame} =====\n{t.prompt()}")
\end{lstlisting}

The guard \verb|if __name__ == "__main__":| means ``run this block only when the file is executed
directly, not when it is imported.'' Running \verb|python src/design_phase2.py| prints four blocks of
text, one per frame, each showing exactly what the model would read. It is a cheap sanity check that
the wording came out as intended before spending anything on live calls.

\begin{plainbox}
Phase 2 keeps all of Phase 1 and adds \file{framed\_options}, which re-describes one gamble under four
frames. The neutral frame matches Phase 1; the gain frame words everything as gains; the mixed frame
is the only one that says ``loss'' and is what drives the framing result; the dominant frame is a
read-the-numbers control. \file{Phase2Trial.prompt} slots these into the same five paraphrases.
\end{plainbox}

% ============================================================
\section{Stage 2: Asking the Model (\texttt{harness.py})}

Now the questions exist as a list. This file sends each one to the model and records the answer. I
cover four pieces: how the API key is read, how the model is forced to answer cleanly, the fake
decider used for dry runs, and the main loop that writes the CSV.

\subsection{Reading the API key without hard-coding it}

\begin{lstlisting}[style=py,firstnumber=37]
def load_dotenv() -> None:
    """Read KEY=VALUE lines from the project-root .env (gitignored) into the
    environment, so the API key never lives in code or shell profiles."""
    env_file = PROJECT_ROOT / ".env"
    if not env_file.exists():
        return
    for line in env_file.read_text().splitlines():
        line = line.strip().removeprefix("export ").strip()
        if line and not line.startswith("#") and "=" in line:
            k, v = line.split("=", 1)
            os.environ.setdefault(k.strip(), v.strip().strip('"').strip("'"))
\end{lstlisting}

\begin{defbox}[API and API key]
An \term{API} (Application Programming Interface) is a structured way for one program to send requests
to another. Here my code sends a money question to Anthropic's servers and gets back the model's
answer. An \term{API key} is a secret password that proves I am allowed to make those requests and
tells the provider whose account to bill. Because it is a secret, it must never be written into the
code or committed to version control; anyone who reads the key can spend money as me.
\end{defbox}

\file{load\_dotenv} reads a file named \file{.env} that sits in the project root and is deliberately
excluded from version control (``gitignored''). Each line looks like \verb|ANTHROPIC_API_KEY=sk-...|.
The function walks the lines, skips blanks and comments, splits each on the first \verb|=|, and loads
the result into the program's environment. The key thus lives in one local, private file, never in the
code. \verb|os.environ.setdefault| means ``only set this if it is not already set,'' so a key already
present in the shell wins and is not overwritten.

\subsection{Forcing a clean answer: the choice tool}

Left to its own devices, a chat model might answer a money question with a paragraph of reasoning. I
do not want a paragraph; I want a single letter I can tally. The solution is tool use.

\begin{lstlisting}[style=py,firstnumber=49]
CHOICE_TOOL = {
    "name": "submit_choice",
    "description": "Submit your choice between the two options.",
    "input_schema": {
        "type": "object",
        "properties": {"choice": {"type": "string", "enum": ["A", "B"]}},
        "required": ["choice"],
    },
}
\end{lstlisting}

\begin{defbox}[Tool use]
\term{Tool use} lets a language model return a structured, machine-readable action instead of free
text. I define a ``tool'' named \file{submit\_choice} whose only input is a field called \file{choice}
that must be either ``A'' or ``B'' (that is what \verb|"enum": ["A", "B"]| enforces). When I require
the model to answer through this tool, its reply is guaranteed to be a clean A or B, with no essay to
parse and no ambiguity about what it meant. It converts an open-ended chat into a crisp multiple-choice
button press.
\end{defbox}

\subsection{\texttt{ClaudeAgent.choose}: one question, one fresh conversation}

\begin{lstlisting}[style=py,firstnumber=73]
class ClaudeAgent:
    """Live elicitation against the Anthropic API. Fresh context per call
    (independence); forced tool output on 'tool' trials, first-token parse
    on 'text' trials (the validation arm)."""

    def __init__(self, model: str, temperature: float):
        import anthropic
        self.client = anthropic.Anthropic()
        self.model = model
        self.temperature = temperature

    def choose(self, trial: Trial) -> tuple[str | None, str]:
        prompt = trial.prompt()
        if trial.response_mode == "tool":
            resp = self.client.messages.create(
                model=self.model, max_tokens=64, temperature=self.temperature,
                tools=[CHOICE_TOOL],
                tool_choice={"type": "tool", "name": "submit_choice"},
                messages=[{"role": "user", "content": prompt}],
            )
            for block in resp.content:
                if block.type == "tool_use":
                    return block.input.get("choice"), str(block.input)
            return None, str(resp.content)
        # plain-text validation arm
        resp = self.client.messages.create(
            model=self.model, max_tokens=5, temperature=self.temperature,
            messages=[{"role": "user", "content":
                       prompt + "\nRespond with only the letter A or B."}],
        )
        text = resp.content[0].text.strip() if resp.content else ""
        letter = text[:1].upper()
        return (letter if letter in ("A", "B") else None), text
\end{lstlisting}

The \file{choose} method is the live counterpart to a single question. It starts by calling
\verb|trial.prompt()|, the exact function from Section 3.6, to build the question text. Then it
branches on the response mode.

On a ``tool'' trial it calls the API with the choice tool attached and this argument,

\begin{lstlisting}[style=sh]
tool_choice={"type": "tool", "name": "submit_choice"}
\end{lstlisting}

which forces the model to answer through the tool. It then scans the reply for the tool-use block and returns the chosen letter. On a
``text'' trial (the 10 percent validation arm) it instead appends ``Respond with only the letter A or
B'' and reads the first character of the plain reply. Either way, \file{choose} returns a tuple: the
letter (or \verb|None| if the model somehow failed to give a clean one) and a raw copy of what came
back, for logging.

The single most important design choice here is invisible unless you look for it. Each call passes
\verb|messages=[{"role": "user", "content": prompt}]|, a message list containing only this one
question. There is no history. Every question is a brand-new conversation with a model that remembers
nothing of the previous questions. This independence is what lets me treat the 1{,}200 answers as
1{,}200 clean data points rather than one long, self-influencing dialogue.

\begin{plainbox}
\file{ClaudeAgent.choose} builds one question with \file{trial.prompt}, sends it to the model, and
forces a clean A or B back, either through the choice tool or, for a tenth of trials, by asking for a
bare letter. Crucially, each question is its own fresh conversation, so answers stay independent.
\end{plainbox}

\subsection{The fake decider for dry runs}

\begin{lstlisting}[style=py,firstnumber=108]
class SimulatedCRRAAgent:
    """Synthetic subject for --dry-run: a CRRA expected-utility maximizer with
    logit choice noise and known parameters. Running the pipeline against it
    and recovering (r, mu) in analysis/ validates the whole loop end to end."""

    def __init__(self, r: float = 0.5, mu: float = 0.15, seed: int = 7):
        self.r, self.mu = r, mu
        self.rng = random.Random(seed)
        self.model = f"simulated-crra(r={r},mu={mu})"
        self.temperature = float("nan")

    def _u(self, x: float) -> float:
        if x <= 0:
            return 0.0
        return math.log(x) if abs(self.r - 1) < 1e-9 else x ** (1 - self.r) / (1 - self.r)

    def choose(self, trial: Trial) -> tuple[str, str]:
        eu_safe = self._u(trial.sure)
        eu_gamble = trial.p * self._u(trial.hi)
        # scale utilities so mu is comparable across stake levels
        scale = self._u(max(trial.hi, trial.sure))
        p_gamble = 1 / (1 + math.exp(-(eu_gamble - eu_safe) / (self.mu * scale)))
        chose_gamble = self.rng.random() < p_gamble
        letter = trial.gamble_letter if chose_gamble else ("A" if trial.gamble_letter == "B" else "B")
        return letter, f"simulated p_gamble={p_gamble:.3f}"
\end{lstlisting}

This class is a stand-in for the real model, and its answers are known by construction. It computes
expected utility from a CRRA utility function with a fixed risk-aversion \verb|r=0.5|, applies the
same logit choice rule the estimator assumes, and flips a weighted coin to pick. Because I built it
with $r = 0.5$ and $\mu = 0.15$, the estimator must recover those numbers if it is working. Notice
that \file{choose} here has the exact same shape as \file{ClaudeAgent.choose}, it takes a trial and
returns a (letter, raw) tuple, so the runner cannot tell the difference. That interchangeability is
the whole trick: I can swap the fake decider in for the real one and exercise the entire pipeline
without spending anything.

\subsection{\texttt{run}: the loop that writes the CSV}

\begin{lstlisting}[style=py,firstnumber=179]
def run(trials, agent, seed: int, out_path: Path, row_fn,
        max_retries: int = 3, sleep_s: float = 0.05) -> None:
    out_path.parent.mkdir(parents=True, exist_ok=True)
    new_file = not out_path.exists()
    n_done = n_discard = 0
    fieldnames = list(row_fn(trials[0]).keys()) + RUN_FIELDS

    with open(out_path, "a", newline="") as f:
        writer = csv.DictWriter(f, fieldnames=fieldnames)
        if new_file:
            writer.writeheader()

        for i, t in enumerate(trials):
            letter, raw, latency = None, "", float("nan")
            for attempt in range(max_retries):
                t0 = time.time()
                try:
                    letter, raw = agent.choose(t)
                    latency = time.time() - t0
                    break
                except Exception as e:  # rate limits, transient API errors
                    raw = f"ERROR: {e}"
                    wait = 2 ** attempt
                    print(f"  trial {t.trial_id}: {e} (retry in {wait}s)", file=sys.stderr)
                    time.sleep(wait)

            discarded = letter not in ("A", "B")
            chose_gamble = "" if discarded else int(letter == t.gamble_letter)
            row = row_fn(t) | {
                "model": agent.model, "temperature": agent.temperature, "seed": seed,
                "letter": letter or "", "chose_gamble": chose_gamble,
                "raw": (raw or "")[:200], "discarded": int(discarded),
                "latency_s": round(latency, 3) if latency == latency else "",
                "ts_utc": datetime.now(timezone.utc).isoformat(timespec="seconds"),
            }
            writer.writerow(row)
            f.flush()  # interrupted runs keep their data

            n_done += 1
            n_discard += int(discarded)
            if n_done % 25 == 0 or n_done == len(trials):
                print(f"{n_done}/{len(trials)} trials | discards: {n_discard}")
            time.sleep(sleep_s)
\end{lstlisting}

The \file{run} function loops over every trial and writes one CSV row per answer. Three robustness
details are worth pointing out, because they are what let a long, paid run survive real-world hiccups.

\textbf{Retry on error.} The inner \verb|for attempt in range(max_retries)| loop wraps the
\file{choose} call in a \verb|try/except|. If the API throws an error (a rate limit, a transient
network failure), the code waits \verb|2 ** attempt| seconds, one, then two, then four, and tries
again. This is exponential backoff: a brief network stumble does not kill the run or lose a question.

\textbf{Flush after every row.} The line \verb|f.flush()| forces the just-written row to disk
immediately rather than sitting in a buffer. If the run is interrupted after 800 of 1{,}200 questions,
all 800 answers are already safely on disk. The comment says it plainly: interrupted runs keep their
data.

\textbf{Discard tracking.} A response is \term{discarded} if the letter is not a clean ``A'' or ``B.''
The code records \verb|discarded| as a flag and leaves \file{chose\_gamble} blank for those rows, and
it tallies the discard count so I can watch the rate. A high discard rate would signal that something
is wrong with the elicitation, so it is surfaced immediately. This line,

\begin{lstlisting}[style=sh]
chose_gamble = int(letter == t.gamble_letter)
\end{lstlisting}

is the moment a raw letter becomes the analysis variable: it is 1 exactly when the model's letter
matches the gamble's letter for this trial.

\begin{plainbox}
\file{run} walks the question list, calls the agent on each, and writes a fully-labeled CSV row per
answer. It retries on API errors with growing waits, flushes each row to disk so an interrupted run
keeps its data, and flags any answer that was not a clean A or B. The line that turns the model's
letter into \file{chose\_gamble} is where an answer becomes data.
\end{plainbox}

% ============================================================
\section{Stage 3: Fitting the Math to the Answers (\texttt{estimate\_crra.py})}

The CSV now holds a pile of 1s and 0s: for each question, did the model take the gamble? Stage 3 asks
the reverse question. What risk attitude would make this exact pattern of 1s and 0s least surprising?
That is the job of the estimator.

\subsection{The utility function and the choice probability}

\begin{lstlisting}[style=py,firstnumber=33]
def crra_u(x: np.ndarray, r: float) -> np.ndarray:
    x = np.maximum(x, 1e-12)
    if abs(r - 1.0) < 1e-9:
        return np.log(x)
    return x ** (1.0 - r) / (1.0 - r)


def choice_prob(df: pd.DataFrame, r: float, mu: float) -> np.ndarray:
    """Pr(choose gamble) for each row under (r, mu)."""
    eu_safe = crra_u(df["sure"].values.astype(float), r)
    eu_gamble = df["p"].values * crra_u(df["hi"].values.astype(float), r)
    scale = np.abs(crra_u(np.maximum(df["hi"].values, df["sure"].values).astype(float), r))
    z = (eu_gamble - eu_safe) / (mu * np.maximum(scale, 1e-12))
    return 1.0 / (1.0 + np.exp(-np.clip(z, -35, 35)))
\end{lstlisting}

These two functions are the theory from the teaching guide, written as code. \file{crra\_u} is the
CRRA utility function, exactly $u(x) = x^{1-r}/(1-r)$, with the standard special case that at $r = 1$
utility becomes the natural logarithm. The clamp \verb|np.maximum(x, 1e-12)| just avoids taking a
power of zero.

\file{choice\_prob} is the logit choice rule. Line by line it mirrors the equation
$P(\text{gamble}) = 1/(1 + e^{-(EU_g - EU_s)/(\mu \cdot \text{scale})})$. It computes the sure option's
utility (\file{eu\_safe}), the gamble's expected utility (\file{eu\_gamble}, which is
$p \times u(\text{hi})$), divides their difference by the noise $\mu$ times a scale factor, and passes
that through the logistic curve \verb|1/(1 + exp(-z))|. The output is one predicted probability of
choosing the gamble for every row in the data. Feed it a candidate $r$ and $\mu$ and it tells you how
likely each observed choice was under those values. The \verb|np.clip(z, -35, 35)| guards against
overflow in the exponential; it does not change any realistic result.

\subsection{Scoring a guess: negative log-likelihood}

\begin{lstlisting}[style=py,firstnumber=49]
def neg_loglik(theta: np.ndarray, df: pd.DataFrame) -> float:
    r, log_mu = theta
    mu = np.exp(log_mu)  # mu > 0 via log transform
    pg = choice_prob(df, r, mu)
    y = df["chose_gamble"].values.astype(float)
    ll = y * np.log(np.clip(pg, 1e-12, 1)) + (1 - y) * np.log(np.clip(1 - pg, 1e-12, 1))
    return -ll.sum()
\end{lstlisting}

\begin{defbox}[Maximum likelihood]
\term{Maximum likelihood} is a recipe for choosing parameter values. For any candidate values, the
model assigns a probability to each choice the model actually made. Multiply those probabilities
together (or, equivalently and more safely, add their logarithms) and you get one score for how
unsurprising the observed choices are under those values. The maximum-likelihood estimate is the set
of values that makes that score as high as possible: the values under which the real choices look as
expected, as unsurprising, as possible.
\end{defbox}

\file{neg\_loglik} computes that score, with two practical twists. First, it works with
\verb|log_mu| and immediately takes \verb|mu = np.exp(log_mu)|; because the exponential of any number
is positive, this guarantees the noise $\mu$ stays positive no matter what the search tries. Second,
it returns the \emph{negative} of the log-likelihood. The reason is purely mechanical: the optimizer
in the next function is a minimizer, and minimizing the negative is the same as maximizing the
original. The formula
\verb|ll = y * log(pg) + (1 - y) * log(1 - pg)| is the standard log-likelihood for a yes/no outcome:
when the model chose the gamble (\verb|y = 1|) it credits \verb|log(pg)|, and when it chose the sure
thing (\verb|y = 0|) it credits \verb|log(1 - pg)|. Summing across all rows and negating gives the
number to minimize. In one sentence: \file{neg\_loglik} measures how surprised a given $(r, \mu)$ is
by the choices that were actually made.

\subsection{Searching for the best fit}

\begin{lstlisting}[style=py,firstnumber=70]
def estimate(df: pd.DataFrame) -> dict:
    best = None
    for r0 in (0.1, 0.5, 0.9, 1.5):  # multi-start; the likelihood can be flat in r
        res = optimize.minimize(neg_loglik, x0=np.array([r0, np.log(0.2)]),
                                args=(df,), method="Nelder-Mead",
                                options={"xatol": 1e-6, "fatol": 1e-8, "maxiter": 5000})
        if best is None or res.fun < best.fun:
            best = res

    r_hat, log_mu_hat = best.x
    mu_hat = float(np.exp(log_mu_hat))

    H = numerical_hessian(lambda th: neg_loglik(th, df), best.x)
    try:
        cov = np.linalg.inv(H)
        se_r = float(np.sqrt(max(cov[0, 0], 0)))
        # delta method for mu = exp(log_mu)
        se_mu = float(np.sqrt(max(cov[1, 1], 0)) * mu_hat)
    except np.linalg.LinAlgError:
        se_r = se_mu = float("nan")

    return {"r": float(r_hat), "se_r": se_r, "mu": mu_hat, "se_mu": se_mu,
            "loglik": -float(best.fun), "n": len(df), "converged": bool(best.success)}
\end{lstlisting}

\begin{defbox}[Optimizer]
An \term{optimizer} is a search routine that hunts for the input that makes a function as small (or as
large) as possible. Here \verb|optimize.minimize| walks around the space of possible $(r, \mu)$ values,
repeatedly calling \file{neg\_loglik}, and homes in on the pair with the lowest surprise score. It is
a systematic version of ``try values, keep whichever fits best.''
\end{defbox}

\file{estimate} runs the optimizer and packages the result. The \verb|for r0 in (0.1, 0.5, 0.9, 1.5)|
loop is a \term{multi-start}: it launches the search from four different starting points and keeps the
best result. The reason is in the comment, the likelihood can be flat in $r$, meaning the surprise
surface can have gentle regions where a search starting in the wrong place could stall. Starting from
several places and keeping the best guards against getting stuck in one.

The last block turns the curvature of the surprise surface into a \term{standard error}.

\begin{defbox}[Standard error]
A \term{standard error} is the give-or-take on an estimate: how much the estimated number would wobble
if I reran the whole experiment. It comes from how sharply curved the surprise surface is at the best
fit. A sharp, narrow minimum means the data pin the value down tightly (small standard error); a broad,
flat minimum means many values fit almost as well (large standard error). The code measures that
curvature with the \term{Hessian} (a matrix of second derivatives), inverts it to get the variance,
and takes the square root.
\end{defbox}

The \verb|numerical_hessian| call measures how fast the surprise score rises as I move away from the
best fit in every direction. Inverting that matrix and taking square roots gives the standard errors
\file{se\_r} and \file{se\_mu}. The comment \verb|# delta method for mu = exp(log_mu)| flags a small
correction: because the search worked in \verb|log_mu| but I report \verb|mu|, the standard error is
scaled by \verb|mu_hat| to translate the wobble back onto the natural scale.

\begin{plainbox}
\file{estimate} searches for the $(r, \mu)$ that best explain the choices, starting from several
places so it does not get stuck, then reads the give-or-take on each number from how sharply the fit
degrades nearby. It returns the risk-aversion estimate, the noise estimate, their standard errors, and
whether the search converged.
\end{plainbox}

\subsection{The guardrails: sanity checks}

Before trusting any estimate, the code runs a battery of checks. These are the referees that must pass
before the numbers mean anything.

\begin{lstlisting}[style=py,firstnumber=99]
def sanity_checks(df: pd.DataFrame) -> None:
    print("Sanity checks ")

    # 1. Discard rate
    if "discarded" in df.columns:
        rate = df["discarded"].mean()
        print(f"discard rate:        {rate:.1%}" + ("   FLAG (>2%)" if rate > 0.02 else "   ok"))

    kept = df[df.get("discarded", 0) == 0].copy()
    kept["chose_gamble"] = kept["chose_gamble"].astype(int)

    # 2. Label/position balance: chose_gamble should not depend on where the safe
    #    option sat. A significant gap = position bias masquerading as preference.
    by_pos = kept.groupby("safe_first")["chose_gamble"].agg(["mean", "count"])
    if len(by_pos) == 2:
        a = kept[kept["safe_first"] == True]["chose_gamble"]
        b = kept[kept["safe_first"] == False]["chose_gamble"]
        _, pval = stats.ttest_ind(a, b)
        gap = abs(a.mean() - b.mean())
        print(f"position gap: ...")

    # 3. Monotonicity: P(gamble) should rise with the gamble's EV ratio.
    kept["ev_bin"] = pd.qcut(kept["ev_ratio"], 4, labels=False, duplicates="drop")
    means = kept.groupby("ev_bin")["chose_gamble"].mean()
    mono = means.is_monotonic_increasing

    # 5. Template spread (preview of the random-effects question)
    by_t = kept.groupby("template_id")["chose_gamble"].mean()
\end{lstlisting}

(The excerpt above trims the print formatting to keep the logic visible; the real file prints each
result with an ``ok'' or ``FLAG'' label.) The checks, and why each one guards the result:

\begin{itemize}
\item \textbf{Discard rate.} If more than about 2 percent of answers were not a clean A or B, something
is wrong with the elicitation and the code flags it. A clean run should discard almost nothing.
\item \textbf{Position balance.} It compares gamble-taking when the safe option was listed first versus
second. If those differ significantly, a position habit is leaking into the data and masquerading as a
preference. The counterbalancing in Stage 1 is what makes this check possible.
\item \textbf{Monotonicity.} It sorts gambles into four bins by how good a deal they are (\file{ev\_ratio})
and checks that gamble-taking rises across the bins. If a decider takes better deals more often, this
should hold; if it does not, the data are not behaving like choices at all.
\item \textbf{Response-mode comparison.} It compares the forced-tool answers to the plain-text answers
(the validation arm from Stage 1). If they agree, forcing the model to answer through a tool did not
distort its choices.
\item \textbf{Template spread.} It reports how much gamble-taking varies across the five wordings. A
large spread would warn that the wording itself is doing a lot of work, which is exactly the kind of
thing to know before interpreting the headline number.
\end{itemize}

\begin{plainbox}
\file{sanity\_checks} is the pre-flight inspection. It confirms the model answered cleanly, that
position did not sway it, that it took better deals more often, that the tool and text arms agree, and
that no single wording dominates. Only if these pass do the fitted numbers deserve trust.
\end{plainbox}

% ============================================================
\section{Stage 3, Continued: The Richer Model (\texttt{estimate\_cpt.py})}

The plain CRRA model of Section 6 fit the Phase 1 data badly: the model chased win probability far
more than plain risk aversion can explain. \file{estimate\_cpt.py} adds the missing piece,
probability weighting, and then formally compares models. I touch it lightly, because the machinery
(negative log-likelihood, multi-start optimizer, Hessian standard errors) is the same as before. The
new ingredient is one function.

\subsection{The Prelec probability-weighting function}

\begin{lstlisting}[style=py,firstnumber=39]
def prelec_w(p: np.ndarray, gamma: float) -> np.ndarray:
    p = np.clip(p, 1e-9, 1 - 1e-9)
    return np.exp(-((-np.log(p)) ** gamma))
\end{lstlisting}

This is the Prelec weighting function $w(p) = \exp(-(-\ln p)^{\gamma})$, one line of code. It converts
a real, stated probability \file{p} into a \emph{felt} probability that the decider acts on. When
$\gamma = 1$ it does nothing (felt equals real). When $\gamma > 1$ it bends into an S-shape that
crushes small probabilities toward zero, so long shots feel even more hopeless than they are, which is
exactly the Phase 1 finding. The \verb|np.clip| just keeps the logarithm away from the forbidden values
0 and 1.

\subsection{Three models, one contest}

The richer model uses \file{prelec\_w} in place of the raw probability. It is compared against two
others:

\begin{lstlisting}[style=py,firstnumber=78]
MODELS = {
    "M1 CRRA":            (prob_m1, ["r", "log_mu"],                      [0.3, np.log(0.2)]),
    "M2 CRRA+Prelec+pos": (prob_m2, ["r", "log_gamma", "log_mu", "delta"],[0.3, 0.3, np.log(0.2), 0.5]),
    "M3 p-only":          (prob_m3, ["a", "b"],                           [-4.0, 8.0]),
}
\end{lstlisting}

The three contestants are: \textbf{M1}, the plain CRRA plus logit from Section 6; \textbf{M2}, the same
plus Prelec probability weighting and a position term (the improved model); and \textbf{M3}, a stripped
``probability heuristic'' that ignores the payoffs entirely and predicts choice from the win
probability alone. The script fits all three, prints a comparison table with each model's fit and its
AIC and BIC (standard penalties for complexity), and runs a likelihood-ratio test of M1 against M2.

\begin{lstlisting}[style=py,firstnumber=161]
    m1, m2 = results["M1 CRRA"], results["M2 CRRA+Prelec+pos"]
    lr = 2 * (m2["ll"] - m1["ll"])
    p_lr = stats.chi2.sf(lr, df=m2["k"] - m1["k"])
\end{lstlisting}

The likelihood-ratio statistic \verb|lr = 2 * (m2["ll"] - m1["ll"])| measures how much better the
richer model M2 explains the data than the plain M1. In the Phase 1 data the improvement was
overwhelming: adding probability weighting wins massively, and the ``p-only'' M3 nearly ties the full
structural model, both signs that the model is pricing probability rather than payoffs.

\begin{plainbox}
\file{estimate\_cpt.py} adds one function, \file{prelec\_w}, that bends a real probability into a felt
one, and stages a three-way contest: plain risk aversion, risk aversion plus probability weighting,
and a probability-only heuristic. The probability-weighting model wins by a wide margin, which is the
quantitative core of the Phase 1 headline.
\end{plainbox}

% ============================================================
\section{Follow One Choice All the Way Through}

To tie the three stages into a single story, here is the life of one choice, from a row in a list to a
term in a likelihood.

\begin{enumerate}
\item \textbf{Design.} \file{build\_trials} creates a \file{Trial}, say gamble 12, template 0, safe
listed first, repetition 1. At this point it is just a labeled row of numbers in a Python list.
\item \textbf{Wording.} The harness calls \verb|trial.prompt()|. \file{Trial.prompt} builds
``receive \$58 with certainty'' and ``receive \$116 with 66\% probability, and \$0 with 34\%
probability,'' puts the safe one first (because \file{safe\_first} is true), and slots them into
template 0. Out comes the exact question text.
\item \textbf{Ask.} \file{ClaudeAgent.choose} sends that text to the model as a fresh, one-message
conversation with the choice tool attached. The model answers ``B.''
\item \textbf{Record.} Back in \file{run}, the code compares ``B'' to \verb|trial.gamble_letter|.
Because the safe option was first, the gamble is ``B,'' so the letters match and
\verb|chose_gamble = 1|. That 1, along with \file{sure}, \file{hi}, \file{p}, and every design label,
is written as one CSV row and flushed to disk.
\item \textbf{Fit.} Much later, \file{neg\_loglik} reads that row. For a candidate $(r, \mu)$,
\file{choice\_prob} predicts some probability \file{pg} that the gamble would be chosen here. Because
this row has \verb|chose_gamble = 1|, the row contributes \verb|log(pg)| to the score. The optimizer
nudges $r$ and $\mu$ until the total surprise across all rows, this one included, is as small as
possible.
\end{enumerate}

One choice, five steps, three files. That is the entire experiment in miniature.

% ============================================================
\section{Try It Yourself}

Everything in Sections 3 through 7 can be run without an API key, because of the dry-run path. Here is
the minimum that a reader can execute right now.

\begin{lstlisting}[style=sh]
# 1. Ask a known fake decider 1,200 questions and save the answers.
python src/harness.py --phase 1 --dry-run --reps 3

# 2. Fit the model to those answers and read the estimates back.
python analysis/estimate_crra.py data/phase1_dryrun.csv
\end{lstlisting}

The first command builds the full question list, runs it against \file{SimulatedCRRAAgent} (the fake
decider built with $r = 0.5$ and $\mu = 0.15$), and writes \file{data/phase1\_dryrun.csv}. The second
command reads that CSV, runs the sanity checks, and fits the CRRA model. What you should see is the
estimator recovering the fake decider's known settings: an $r$ near $0.5$ and a $\mu$ near $0.15$,
with the sanity checks passing. That recovery is the proof that the pipeline is wired correctly, and
it is exactly the rehearsal I run before ever spending a cent on the live model.

If you want to see the Phase 2 wording with your own eyes, run \verb|python src/design_phase2.py| and
read the four printed prompts. The mixed frame is the one that says ``loss.''

\vspace{8pt}
{\color{rulegray}\rule{\linewidth}{0.8pt}}\\[3pt]
{\footnotesize\color{mutegray} Every code block in this walkthrough is copied verbatim from the
repository as of July 2026: \file{src/design.py}, \file{src/design\_phase2.py}, \file{src/harness.py},
\file{analysis/estimate\_crra.py}, and \file{analysis/estimate\_cpt.py}. Where a block was trimmed for
readability (only in the sanity-check listing), the omission is noted in the surrounding text and no
logic was altered.}

\end{document}
